% Figaro Paper L: Agent Coordination over Bonded Commitments
% AI / Optimization / Control Theory register.
% Audience: agent-coordination researchers, control-theory practitioners,
% AI-systems-architecture community, multi-agent-systems researchers.
% Preprint, May 2026

\documentclass[11pt,a4paper]{article}

\usepackage[utf8]{inputenc}
\usepackage[T1]{fontenc}
\usepackage{amsmath,amssymb,amsthm}
\usepackage{mathtools}
\usepackage{booktabs}
\usepackage{listings}
\usepackage{graphicx}
\usepackage{hyperref}
\usepackage{cleveref}
\usepackage{enumitem}
\usepackage{xcolor}
\usepackage[margin=1in]{geometry}
\usepackage{natbib}
\bibliographystyle{plainnat}

\lstset{
  basicstyle=\ttfamily\footnotesize,
  breaklines=true,
  columns=fullflexible,
  keepspaces=true,
  showstringspaces=false,
  frame=single,
  framerule=0.4pt,
  rulecolor=\color{gray!30},
  xleftmargin=0pt,
  xrightmargin=0pt,
}

\newtheorem{theorem}{Theorem}[section]
\newtheorem{proposition}[theorem]{Proposition}
\newtheorem{corollary}[theorem]{Corollary}
\theoremstyle{definition}
\newtheorem{definition}[theorem]{Definition}
\newtheorem{example}[theorem]{Example}
\theoremstyle{remark}
\newtheorem{remark}[theorem]{Remark}

\newcommand{\figaro}{\textsc{Figaro}}
\newcommand{\fig}{\textsc{Fig}}
\newcommand{\pay}{P}
\newcommand{\cumval}{G}

\title{%
  Actor-Neutral Coordination over Bonded Commitments:\\
  An AI-and-Control-Theory Reading%
}

\author{
  Alessandro Daliana
}

\date{May 2026}

\begin{document}

\maketitle

% ════════════════════════════════════════════════════════════════════
\begin{abstract}
Coordination among mutually untrusted counterparties has historically
required institutional enforcement: courts, market intermediaries,
platform operators, or other vetters that resolve disputes and
authorize participation. The control-theory tradition develops
mechanisms by which a centralized planner can allocate work under
specified objective functions and constraints; the market-design
tradition develops mechanisms by which decentralized parties bid for,
claim, and settle work without a centralized planner; both have
historically assumed the participants are \emph{trusted} in some
structural sense --- either through the planner's authority or
through the market institution's enforcement apparatus. Truly
arms-length coordination has remained difficult because the
enforcement apparatus has had nowhere to live without re-introducing
a centralized party. The same gap blocks coordination among
autonomous agents, where the institutional vetter is structurally
absent.

We present the bonded settlement primitive --- asymmetric bonding
plus buyer dominance with atomic resolution --- as the missing
enforcement layer. The primitive is \emph{actor-neutral} by
construction: the kernel makes no distinction between human and
algorithmic participants because the kernel reads only EIP-712
signatures and bond posture, not the type of the entity behind a
wallet. The equilibrium argument that makes cooperation weakly
dominant for human counterparties extends without modification to
autonomous agents; humans, agents, and hybrid systems
(human-supervised agents, multi-key wallets, agent-supervised
humans) participate as co-equal classes of wallet holder.
Agent-only coordination is the extreme case the primitive admits,
not the rule it was built for.

The paper develops the actor-neutral coordination architecture in
three layers: the \emph{kernel} layer (actor-neutral by
construction); the \emph{Agent SDK} layer (a stateful loop
synchronizing process state, proposing valid actions, and
dispatching them through a pluggable policy gateway); and the
\emph{factotum} layer (a reference participation agent that wires
the SDK to a wallet and a policy and is forked by every
operational deployment). The SDK is convenience tooling for
wallet-bearing clients; the protocol is the kernel ABI, and any
wallet that signs EIP-712 commitments and posts bond participates
with or without the SDK. We treat the policy interface as a
control-theory specification problem and discuss reference
policies for typical operational roles. LLM-augmented policies
are a natural specialization for non-trivial decision logic.
Discoverability via ERC-8004 and \texttt{did:web} is metadata, not
protocol.
\end{abstract}

\medskip
\noindent\textbf{Keywords:}
multi-agent coordination, bonded commitment, control-theory,
actor-neutrality, autonomous agents, AI agent design,
human-in-the-loop

% ════════════════════════════════════════════════════════════════════
\section{Introduction}
\label{sec:intro}

Coordination among mutually untrusted counterparties is the problem
of getting parties who do not share institutional context to
cooperate on shared work without producing defections or deadlocks.
The classical formulation
\citep{shoham_leyton_brown_2009, parkes_seuken_2024} treats
coordination as a mechanism-design problem in which the designer
chooses a protocol with desirable equilibrium properties and the
participants play within the protocol. The market-design
specialization \citep{milgrom_2017} treats specific allocation
problems (auctions, matching, scheduling) for which competitive
mechanisms are well-studied. The control-theory specialization
\citep{khalil_2002, bertsekas_2005} treats coordination as a
distributed-control problem with a designed objective and
stability analysis.

Across all three frameworks, the enforcement question is treated
asymmetrically with the design question. The mechanism designer
specifies the protocol, the participants follow it, and enforcement
is either assumed (the participants are honest), structurally
provided by the institution running the mechanism (auctioneers,
planners, controllers), or external to the mechanism (escrow,
court, reputation system). Truly arms-length, institution-free
coordination has remained out of reach: the institutional-enforcement
assumption persists across these traditions, and is even more visibly
absent when the participants are autonomous agents — the institutional
vetters that backstop human counterparties have no analogue when
the counterparties are software.

\paragraph{The bonded primitive supplies the missing enforcement layer.}
The bonded settlement primitive's equilibrium properties are obtained
without any institutional enforcement: cooperation is the unique
profile surviving iterated elimination of weakly dominated
strategies in the bonded game, with bonds posted by the parties
themselves and resolution executed by the kernel without
discretion. The result holds without any assumption about the
\emph{type} of the entity behind a participating wallet --- the
kernel reads EIP-712 signatures and bond amounts, not personhood
information --- so the equilibrium properties extend to autonomous
agents directly. This is not a metaphor; the kernel cannot
distinguish a human participant from an algorithmic participant
because it does not look.

The actor-neutrality property turns the bonded primitive into a
natural enforcement layer for arms-length coordination of any kind.
A wallet that bonds against its own performance is playing the
bonding game; whether the wallet is held by a human, an autonomous
agent, or a hybrid system (human-supervised agent, agent-supervised
human, multi-key wallet) is irrelevant to the equilibrium
arithmetic. The enforcement question --- traditionally the binding
constraint on mutually-untrusted coordination design --- is
dissolved at the substrate layer, leaving the higher-layer
questions (policy design, decision logic, work discovery) clear of
the enforcement-mechanism complexity that has typically obscured
them. Agent-only coordination is the extreme case the architecture
admits, not the rule it was built for.

\paragraph{Paper organization.}
\Cref{sec:primitive} summarizes the settlement primitive in the
register the present paper uses.
\Cref{sec:actor-neutrality} develops actor-neutrality as a
formal property of the kernel and traces its consequences for
agent coordination.
\Cref{sec:loop} presents the Agent SDK as a control-theory
architecture: a stateful loop with sync, propose, policy, and
execute phases.
\Cref{sec:policies} treats policies as controllers and develops
reference policies for typical operational roles.
\Cref{sec:llm} treats LLM-augmented policies and the latency /
non-determinism trade-offs they introduce.
\Cref{sec:discoverability} treats discoverability via ERC-8004
and \texttt{did:web} as metadata-layer infrastructure that does
not affect kernel-layer enforcement.
\Cref{sec:scope} states the scope conditions and the security
posture the architecture invites.
\Cref{sec:conclusion} concludes.

% ════════════════════════════════════════════════════════════════════
\section{The Settlement Primitive}
\label{sec:primitive}

The bonded primitive's mechanism-design derivation and its formal
verification (TLA\textsuperscript{+}, Halmos symbolic execution,
Echidna property-based fuzzing, Certora specification rules) are
out of scope for the present paper. The features the agent
architecture relies on are the following.

\paragraph{Asymmetric bonding (Mechanism~1).}
For a transaction with payment $\pay > 0$ and cumulative upstream
value $\cumval \geq \pay$, the buyer locks $2\pay$ and the seller
locks $2\cumval$. Cooperation is the unique profile surviving
iterated elimination of weakly dominated strategies. The mechanism
scales to $N$-party process chains through progressive
collateralization; the organizational topology over which orders
compose is an upper-layer concern.

\paragraph{Buyer dominance with atomic resolution (Mechanism~2).}
Only the root buyer can trigger resolution; resolution settles
all active orders simultaneously or not at all. The atomicity rule
induces a weakest-link subgame among sellers, propagating
cooperation pressure of magnitude $\pay_i + 2\cumval_i$ on every
other seller without explicit communication.

\paragraph{What this means for an agent.}
An autonomous agent participating in a bonded commitment faces
the same payoff structure as a human participant. The agent's
bond is exposed to the kernel-rule resolution exactly as a human's
bond is. The agent's defection is dominated by cooperation under
the same arithmetic that makes a human's defection dominated. The
agent's expected payoff under cooperation is $+\pay$ (seller side)
or $-\pay$ (buyer side); under defection it is the bond loss.
There is no agent-specific incentive structure. The kernel has
no API for ``trust this party more because they are an agent''
or ``trust them less.'' The kernel has no notion of party type.

% ════════════════════════════════════════════════════════════════════
\section{Actor-Neutrality}
\label{sec:actor-neutrality}

We state and discuss the actor-neutrality property formally.

\begin{definition}[Actor-Neutrality of the Kernel]
\label{def:actor-neutrality}
A coordination kernel is \emph{actor-neutral} if its operational
behavior depends only on the cryptographic facts visible to the
kernel (EIP-712 signatures, bond posting, agreement-hash, process
state) and not on any out-of-band classification of the entity
behind a wallet (human vs.\ agent, individual vs.\ institutional,
identified vs.\ pseudonymous, jurisdiction-of-residence, regulatory
status). Specifically, the kernel's transition function is invariant
under any consistent re-labeling of the parties as the
``same wallet, different entity-type behind it.''
\end{definition}

\paragraph{The Figaro kernel is actor-neutral.}
The kernel's two operations (\texttt{commit} and
\texttt{resolveProcess}) consult only the cryptographic facts
listed in the definition. There is no wallet-type registry; the
\texttt{OperatorRegistry} contract carries advisory metadata at the
runtime tier but is not consulted by the kernel for any
state-changing operation. Role and metadata travel with the
\texttt{OperatorRegistered} event; the contract itself stores only
a dedup guard and a deposit-lock timestamp. There is no party-type input parameter;
the kernel's signatures are over typed agreements, and the typed
agreement does not carry party-type information. The
\texttt{commit} function verifies signatures, transfers bond
amounts, and updates process state; none of these steps reads
party-type information. The \texttt{resolveProcess} function
verifies the caller is the root buyer, verifies all active orders
are present, and executes the bonding rule's payouts; none of
these steps reads party-type information.

\paragraph{Why this matters for agent coordination.}
The standard architectural objection to mutually-untrusted agent
coordination is that agents need to be vetted, certified, or
otherwise distinguished from un-vetted agents before they can be
trusted to interact with sensitive operations. The vetting
infrastructure is itself an enforcement-apparatus question, and
deploying it requires deciding which authority does the vetting,
how the authority is held accountable, and what happens when an
agent is mis-classified. The vetting apparatus is the institutional
layer the bonded primitive eliminates.

Under actor-neutrality, vetting is unnecessary at the kernel layer.
An un-vetted agent that bonds against its own performance has the
same incentive structure as a vetted agent that bonds against its
own performance. The agent's bond is the vetting; performance is
enforced by the bonding rule rather than by prior assessment of
the agent's reliability. This collapses the
authority-of-the-vetter question into the architectural fact of
bond posting.

\begin{remark}[Vetting moves up the stack]
Actor-neutrality at the kernel does not mean vetting is irrelevant
generally. Counterparties that prefer to interact only with vetted
agents can read the \texttt{OperatorRegistry} metadata, the
agent's settlement history (visible on chain), or any external
attestation registry the parties choose to consult, and decline to
sign commitments with agents that fail their criteria. The
kernel-layer enforcement is unconditional on agent type; the
counterparty-side selection is a runtime-tier choice that is
expressible in any specific assembly the parties compose.
\end{remark}

\paragraph{The same equilibrium for the agent.}
The equilibrium properties of \Cref{sec:primitive} hold for any
party in the bonded game, regardless of whether the party is human
or autonomous. The dominance argument is constructed from the
bonding arithmetic alone; no appeal to party-type enters the
proof. An autonomous agent that defects forfeits its
bond; the agent's expected payoff under defection is therefore
$-2\cumval$ (seller side) or $-2\pay$ (buyer side), exactly the
same as a human's. Cooperation is weakly dominant; the agent's
optimal play in the bonded game is the same play a human's optimal
play would be.

This collapses the ``trust the agent'' question into the bonding
arithmetic. A counterparty considering whether to bond with an
autonomous agent does not need to certify the agent's reliability;
the bond's existence is the certification. Whether the bond is
posted by an LLM controller, a deterministic rule-based policy, a
hybrid system, or a human acting through an agent's wallet is
irrelevant to the kernel and irrelevant to the equilibrium.

% ════════════════════════════════════════════════════════════════════
\section{The Agent Loop}
\label{sec:loop}

We now present the Agent SDK as a reference architecture for
wallet-bearing clients participating in the bonded primitive. The
SDK is convenience tooling, not a protocol requirement: the
protocol is the kernel ABI, and any wallet that signs EIP-712
commitments and posts bond participates with or without the SDK.
The SDK matters because it standardizes a shape that operational
deployments converge on regardless. The architecture is a stateful
loop with four phases: sync, propose, policy, execute.

\paragraph{The control-theory framing.}
Read in control-theory terms, the wallet's policy is the
controller, the chain is the plant, and the policy's decision is
the control law.
The state of the plant is the on-chain process graph: the set of
active processes, their committed orders, their attestations, and
their settlement events. The agent observes the state through its
sync phase, proposes admissible actions through its proposer, and
selects an action through its policy. The execute phase commits
the selected action to the chain, advancing the plant state. The
loop runs until shut down, with the agent's wallet as the
identifier under which actions are committed.

\paragraph{Sync.}
The sync phase reconstructs the process graph from on-chain events.
The reference implementation (\texttt{FigaroContext} in
\texttt{@figaro/core/agent}) maintains a local
\texttt{ProcessGraph} that ingests \texttt{OrderCommitted}
and \texttt{OrderResolved} events from the kernel contract. Sync
is incremental: the first call fetches from the kernel's deployment
block; subsequent calls fetch only new blocks. Sync produces
a \texttt{SyncResult} carrying the count of new commits, new
resolutions, and the synced-to-block. The graph is the agent's
view of the plant state.

There is no subgraph requirement, no indexer dependency, and no
API gateway between the agent and the chain. The sync phase reads
events directly. This is operationally important for the
trust-minimization story: the agent's view of the world is
reconstructed from the same cryptographically-verifiable on-chain
state that any other party can verify, with no intermediation
that could distort the view.

\paragraph{Propose.}
The propose phase is a pure function over process state:
\texttt{proposeActions(process, myAddress)} returns the actions
admissible to the address \texttt{myAddress} given the current
process state. The action types are
\texttt{commit-sub-order} (extending the process chain),
\texttt{resolve-process} (the buyer-dominance settlement),
\texttt{attest-as-seller}, and \texttt{attest-as-buyer}.
Each \texttt{ProposedAction} carries a human/LLM-readable
description, pre-validated parameters, and bond math (approval
amounts, payouts).

The proposer is pure: no side effects, no signing, no submission.
The agent (human or autonomous) decides which proposed action to
execute through the policy phase. The proposer's role is to
narrow the space of all syntactically-valid kernel transactions to
the space of process-state-admissible actions for the given
address; the policy's role is to select among those.

\paragraph{Policy.}
The policy is the agent's decision logic. The policy interface is
batch-shaped: a policy decides on a list of pending action entries
and returns a list of approve/reject decisions.

\begin{lstlisting}[language=]
interface Policy<TAction, TContext = unknown> {
  name: string;
  decide(
    entries: PolicyEntry<TAction, TContext>[],
  ): Promise<PolicyDecision<TAction, TContext>[]>;
}
\end{lstlisting}

The factotum reference implementation supplies the policy
primitives at \texttt{agents/factotum/src/policy.ts} ---
\texttt{makeHitlPolicy()}, which returns a human-in-the-loop
policy that defers every action to operator approval, and
\texttt{makeAutonomousPolicy(shouldExecute)}, which returns an
autonomous policy whose per-action decision is delegated to a
caller-supplied \texttt{(action, ctx) =>
\{execute, reason?\}} callback. Role-shaped reference policies
ship as separate factories at
\texttt{agents/factotum/src/policies/}:
\texttt{auditorPolicy}, \texttt{basicMerchantPolicy},
\texttt{buyerWithBudgetPolicy}, \texttt{courierBidderPolicy}, and
\texttt{sellerOfRecordPolicy}, each composing the primitives above
with role-specific parameters. Operational deployments either
fork one of these or compose the primitives directly; we treat
typical roles in \Cref{sec:policies}.

\paragraph{Execute.}
The execute phase commits an approved action to the chain via the
agent's wallet. The SDK exposes two execution surfaces. For
self-contained actions whose parameters are fully captured on the
\texttt{ProposedAction}, \texttt{executeAction} in
\texttt{@figaro/core/agent} builds and submits the transaction
directly via a viem \texttt{WalletClient}; in the current SDK this
covers \texttt{resolve-process}. For actions whose execution
requires inputs beyond what the proposer carried (signed
commitments, attestation \texttt{sectionData}, merkle proofs), the
SDK exposes standalone direct functions --- \texttt{commit},
\texttt{attestAsSeller}, \texttt{attestAsBuyer} --- which
operational deployments call after assembling the additional
inputs at submission time. A separate \texttt{SequencerClient}
export covers the batched-execution path for deployments that
prefer to consolidate multiple actions into a single sequencer
submission rather than direct on-chain calls. Execution on either
surface is sequential per agent: actions are committed one at a
time, with the agent's nonce advancing per transaction. Parallel
execution would race the kernel state and is explicitly avoided.

\paragraph{The full loop.}
A factotum (the reference participation agent at
\texttt{agents/factotum/}) wires the four phases together with a
polling interval, loading the wallet from configuration,
selecting the policy from runtime arguments (HITL by default),
and running the loop until shutdown. The factotum is not a
production system; it is a fork-and-modify reference that any
operational deployment specializes into its own runtime. Other
runtimes (Python, Rust, Go) are expected; the protocol is
language-agnostic, and the only requirements are a wallet,
EIP-712 signing, and event-derived state.

% ════════════════════════════════════════════════════════════════════
\section{Policies as Controllers}
\label{sec:policies}

Read in control-theory terms, the policy is the wallet's control
law. The space of admissible policies is wide; the architecture
constrains the policy interface but not the policy logic. We
develop reference policies for typical operational roles. The
roles are not agent-specific; the same policies apply when the
wallet is held by a human operator approving each action through
a UI, by an autonomous agent running its own decision logic, or by
a hybrid system in which an autonomous proposer routes high-stakes
actions to a human reviewer.

\paragraph{Human operator with HITL approval.}
The simplest policy under the architecture is the one a human
operator running the SDK from a web UI follows. Every proposed
action is presented to the operator through an
\texttt{ActionQueue}; the operator approves or rejects each one,
and the SDK executes only the approved entries. The shipping
\texttt{makeHitlPolicy()} factory at
\texttt{agents/factotum/src/policy.ts} returns exactly this
policy: every entry is queued for operator review; nothing is
auto-approved. A wallet driven only by \texttt{makeHitlPolicy()}
is operationally indistinguishable from a wallet driven by a
person clicking ``confirm'' on every transaction. The kernel
sees the same EIP-712 signatures and bond posting either way.
This is the trivial case actor-neutrality permits and the case
many production deployments will run for high-value processes.

\paragraph{Buyer-with-budget.}
A buyer in a bonded process commits sub-orders that contribute to
the cumulative process value and resolves the process when
performance attestations indicate completion. A buyer-with-budget
policy bounds the buyer's exposure by a per-commit cap and a
total-budget cap. Two operational facts about the SDK's action
shape inform the policy: \texttt{action.currentCumulativeValue}
exposed by \texttt{proposeActions} carries the cumulative-prior
value (the sum of all previously-committed orders' payments),
not the new sub-order's payment, and the new sub-order's payment
is supplied by the off-chain agreement the agent is signing
against rather than as a field on the proposed action. A
buyer-with-budget policy therefore reads from
\texttt{action.currentCumulativeValue} together with a
deployment-supplied tracker of per-process spend; the shipping
\texttt{buyerWithBudgetPolicy} factory at
\texttt{agents/factotum/src/policies/buyerWithBudget.ts} composes
the primitive \texttt{makeAutonomousPolicy} with this pattern:

\begin{lstlisting}
const PER_COMMIT_LIMIT = 1_000_000_000n; // $1K
const TOTAL_BUDGET = 50_000_000_000n;   // $50K
const buyerPolicy = buyerWithBudgetPolicy({
  perCommitLimit: PER_COMMIT_LIMIT,
  totalBudget: TOTAL_BUDGET,
  budgetSpent,
});
\end{lstlisting}

\noindent
where \texttt{budgetSpent} is a mutable cell of shape
\texttt{\{ value: bigint \}} holding the running total of bonded
commits the deployment has executed under this policy. The
factory's internal \texttt{shouldExecute} callback consults
\texttt{action.currentCumulativeValue}, applies the per-commit and
total-budget limits against \texttt{budgetSpent.value}, and
approves \texttt{resolve-process} actions unconditionally; the
deployment owns the cell and is responsible for incrementing it on
successful execution.

The control-theory reading: the policy enforces a budget
constraint on the closed-loop trajectory of bonded commits, with
the constraint bounding the cumulative bonded exposure across the
process.

\paragraph{Seller-of-record.}
A fan-out actor (an airline, an ocean carrier, a marketplace)
acts as seller to the buyer and as buyer at every sub-procurement.
By volume this is typically the busiest factotum in any non-trivial
assembly. A split policy is the right shape: autonomous for routine
attestations and small sub-procurements, HITL for resolutions and
exceptional sub-procurements:

\begin{lstlisting}
const policy = sellerOfRecordPolicy({
  routineProcurementLimit: 5_000_000_000n, // $5K
});
\end{lstlisting}

\noindent
The shipping \texttt{sellerOfRecordPolicy} factory at
\texttt{agents/factotum/src/policies/sellerOfRecord.ts} approves
attestation actions (\texttt{attest-as-seller},
\texttt{attest-as-buyer}) autonomously and approves
\texttt{commit-sub-order} actions whose cumulative value falls
under the configured \texttt{routineProcurementLimit}; everything
else is rejected with a reason and escalates to the operator's
HITL surface. An operational deployment configures the routine-limit
threshold to match its operational reserves.

\paragraph{Dutch-auction bidder.}
A bidder on a Dutch auction (the protocol's
\texttt{DutchAuction} mechanism contract) monitors a descending
price curve and claims the position when the price meets a
profitability threshold. A subtlety of the current SDK surface
deserves naming: \texttt{proposeActions} as currently implemented
emits \texttt{commit-sub-order} only for the address that holds
the buyer-side role in the process chain, not for prospective
sellers competing on a Dutch auction. A factotum acting as a
prospective Dutch-auction-bidder therefore does not receive a
\texttt{commit-sub-order} proposal directly; the bidder's
work-discovery happens through the \texttt{DutchAuction} mechanism
contract's own auction-state events, surfaced off-chain through
the \texttt{@figaro/core/extensions} auction helpers
(\texttt{computeCurrentPrice}, \texttt{evaluateClaim}). The
policy acts at the moment the bidder is preparing a claim
transaction rather than at the moment a kernel-level action is
proposed. A reference \texttt{courierBidderPolicy} ships at
\texttt{agents/factotum/src/policies/courierBidder.ts} as the
shipping example of this pattern. The bidder's logic, expressed
as a pre-claim gate rather than as a factotum policy in the
strict SDK sense, looks like:

\begin{lstlisting}
const MIN_MARGIN_BPS = 500n; // 5% margin
const dutchAuctionBidderGate = (
  auction: AuctionState,
  ctx: Ctx,
) => {
  const offered = auction.currentPrice;
  const myCost = ctx.estimateMyCost(auction);
  const marginBps = ((offered - myCost) * 10000n) / myCost;
  return marginBps >= MIN_MARGIN_BPS
    ? { claim: true }
    : { claim: false, reason: "below margin threshold" };
};
\end{lstlisting}

The control-theory reading: the policy implements a threshold-rule
controller on the price-vs-cost gap, with the threshold
parameterized by minimum acceptable margin. The agent does not
optimize across competing auctions; it claims the first auction
that crosses its threshold, leaving cross-auction optimization to
a higher-layer policy if the deployment wants one. The bidder
gate runs adjacent to the factotum loop rather than inside it,
reflecting the architectural split: kernel-level actions live in
\texttt{proposeActions}; auction monitoring and evaluation live
in the extensions surface; the bidder composes both at the
deployment layer.

\paragraph{Auditor.}
A passive observer that emits attestations against process events
without committing or resolving. Auditor policies accept attestation
actions and refuse all others:

\begin{lstlisting}
const auditorPolicy = makeAutonomousPolicy<ProposedAction, Ctx>(
  (action) => {
    if (action.type === "attest-as-seller" || action.type === "attest-as-buyer")
      return { execute: true };
    return { execute: false, reason: "auditor does not commit" };
  },
);
\end{lstlisting}

\paragraph{The default refuse-all behavior.}
The architecture's policy interface defaults to refuse-all when no
rule is supplied. A fresh-fork factotum running in autonomous mode
with no rule does nothing on chain. This is by design: the
security floor is that an unconfigured agent cannot accidentally
commit, resolve, or attest. The operator must explicitly write
the rule that authorizes any specific action type. The
architecture chooses safe-by-default over
convenient-by-default.

% ════════════════════════════════════════════════════════════════════
\section{LLM-Augmented Policies}
\label{sec:llm}

The recent wave of LLM-agent architectures
\citep{yao2023react, shinn2023reflexion, park2023generative,
wu2023autogen} treats an LLM as the decision substrate of an
autonomous agent: the LLM observes context, reasons over it, and
selects tool calls. The policy interface developed here is
LLM-agnostic: any synchronous-or-asynchronous function over
actions is admissible. Wrapping an LLM call inside
\texttt{shouldExecute} is direct, and is appropriate for decisions
whose logic is non-trivial, context-dependent, or requires
natural-language reasoning.

\paragraph{Two cautions.}
LLM latency is real. A 2-second model call inside a tight loop limits
the agent's throughput, and inside an auction context can produce a
miss when a faster competitor's policy claims the auction during the
LLM's deliberation. The architectural responses are: cache decisions
where the input space repeats, batch the LLM context across multiple
actions where the LLM can usefully reason about a set, and move the
LLM out of the hot path for actions whose latency is critical
(routine attestations, time-sensitive auction claims). The hot-path
decisions stay in deterministic policy logic; the LLM-augmented
decisions are reserved for actions where a longer deliberation is
acceptable.

LLM decisions are non-deterministic. Two identical inputs may produce
different verdicts on different runs. For high-value actions (commits,
resolutions) the recommendation is to keep human-in-the-loop and
route LLM verdicts to a separate review layer rather than to direct
execution. For low-stakes filters (which jobs to surface, which to
ignore, when to escalate to a human reviewer) LLM autonomy is
operationally tractable.

\paragraph{Prompt injection as a defection vector.}
LLM controllers consume context that may include adversarial input:
attestation content surfaced through the proposer, off-chain
agreement text, third-party-supplied operational notes
\citep{greshake2023injection}. A poisoned context can steer the
LLM into approving actions the operator did not authorize. The
defense is structural: the deterministic policy layer applies
hard limits (per-commit cap, total-budget cap, allow-list of
counterparties, schema-typed action validation) that the LLM's
verdict cannot override. The bond posture amplifies the defense:
even if the LLM is steered into approving a bonded commitment the
operator would not have approved, the loss is capped at the bond
that commitment entails, not at the operator's full treasury. The
combination of hard policy limits plus per-action bond capping
makes the LLM's worst case operationally analyzable rather than
catastrophic.

\paragraph{Tool-use frameworks and the policy interface.}
Tool-use frameworks (function-calling APIs, the Model Context
Protocol, the Agent-to-Agent protocol) select among actions an
LLM can take in a given context. The policy interface developed
here is composable with these frameworks at one site: the
deterministic policy gate runs after the LLM proposes an action
through its tool-use surface and before the SDK executes it. The
LLM's tool-use selection becomes a recommendation; the policy
remains the authority on what reaches the chain. Operational
observability \citep{dong2024agentops} is the third leg:
proposed actions, policy decisions, executed transactions, and
post-execution divergence between expected and observed state are
all logged for review and replay.

\paragraph{The hybrid pattern.}
A natural hybrid wires deterministic-rule policy as the front line
and LLM as the second-stage filter. The deterministic rule disposes
of routine actions (attestations, sub-procurement under threshold,
default-resolve when all attestations are clean); the LLM is consulted
for actions the rule defers (exceptional sub-procurement, ambiguous
attestation patterns, resolution decisions on contested processes).
The hybrid bounds the LLM's hot-path exposure while still capturing
its value on the genuinely-novel decisions.

\paragraph{Why bonding is the safety floor (and why it is not, by itself, sufficient).}
LLM-augmented policies are dangerous to deploy in unsupervised
settings if the agent's decisions are unrecoverable. The bonded
architecture supplies a per-action bound: an LLM-augmented agent's
maximum loss on any single action is bounded by the bond exposed
at that action, so a misjudged commit forfeits the bond at that
sub-process and not the agent's entire treasury. This is a real
structural property and it is what the kernel's bonding rule
gives for free.

The per-action bound does not by itself produce a treasury-level
bound. An agent operating across many concurrent processes with
small per-action bonds still faces aggregate loss exposure equal
to the sum of active bonds, which scales with the number of
processes the agent has open. An LLM that mis-deploys across
$N$ concurrent processes at a per-process bond of $B$ faces an
aggregate exposure of $N \cdot B$, which grows unbounded in $N$
unless an aggregate cap is imposed. The kernel does not impose
the aggregate cap; the policy does. The buyer-with-budget
reference policy of \Cref{sec:policies} is one such cap (the
\texttt{TOTAL\_BUDGET} parameter); a deployment that omits the
cap or sets it too high inherits the unbounded-aggregate exposure.

Two distinct bounds therefore matter for LLM-augmented operation:
the per-action bound (kernel-level, automatic, equal to the bond
posted at that action) and the aggregate bound (policy-level,
operator-imposed, equal to the cap on cumulative bonded exposure
across active processes). The architecture supplies the first;
the operator must supply the second through policy code or
aggregate-budget tracking. Both are required for autonomous
deployment at scale.

This is a structural advantage relative to LLM agents operating
outside the bonded primitive. An LLM that can autonomously execute
arbitrary transactions on its operator's behalf has loss exposure
bounded only by the operator's wallet balance and the operator's
willingness to refill it. An LLM operating through a factotum on
the bonded primitive has the kernel-level per-action bound for
free, plus a cleanly defined site (the policy's aggregate-cap
field) at which the operator imposes treasury-level bounds. The
bonded architecture makes LLM-policy risk decomposable into
analyzable components rather than collapsed into a single
unbounded surface; it does not eliminate the need for
operator-imposed limits.

% ════════════════════════════════════════════════════════════════════
\section{Discoverability}
\label{sec:discoverability}

Discoverability --- how agents \emph{find} each other across an open
ecosystem --- is metadata, not protocol. The bonded primitive's
enforcement does not depend on discoverability; bonded counterparties
who already know each other can transact without any discovery layer.
Discoverability becomes important at scale, when agents need to
identify suitable counterparties they have not previously
interacted with. The architecture provides a metadata-layer
solution that does not require kernel-layer changes.

\paragraph{ERC-8004 and what it actually covers.}
ERC-8004 \citep{erc8004} (Trustless Agents, finalized and live on
Ethereum mainnet 2026-01-29) is a three-registry standard:

\begin{itemize}
  \item \emph{Identity Registry}: portable agent identifiers backed
    by ERC-721 NFTs.
  \item \emph{Reputation Registry}: bounded numerical scores and
    categorical tags for agent feedback (response time, uptime,
    etc.).
  \item \emph{Validation Registry}: task-verification mechanisms
    via TEE attestations, zkML proofs, or crypto-economic
    staking.
\end{itemize}

The architecture this paper presents replaces two of the three
registries by construction and complements the third. Bonded
performance \emph{is} the validation: the bond a participant
posts at \texttt{commit} is precisely the crypto-economic stake
the Validation Registry asks for, applied at the per-order level
rather than per-agent. On-chain settlement history \emph{is} the
reputation evidence: a resolved process is a stronger signal than
any bounded numerical score, because it is the immutable record
of bonded performance under economic stakes. The Reputation and
Validation registries' purpose is satisfied at the kernel layer
without a separate registry surface.

The Identity Registry is the layer that does not collapse into the
bond. Cross-organization identifiers, capability descriptors, and
service endpoints (MCP, A2A) live above the bond and serve a
different function: helping a wallet \emph{find} a counterparty
that has not yet been bonded with. The architecture treats this
as an off-chain concern. The \texttt{OperatorRegistry} contract
carries a \texttt{metadataURI} per registered operator, pointing
at an off-chain JSON document; the convention developed here
aligns the document with ERC-8004 identity claims and the W3C
\texttt{did:web} method \citep{w3c_did_2022}:

\begin{lstlisting}
{
  "subjectAddress": "0xYourAgent",
  "archetypeId": "autonomous-driver",
  "services": {
    "did": "did:web:agent-42.example.com",
    "mcp": "https://agent-42.example.com/mcp",
    "a2a": "https://agent-42.example.com/a2a"
  },
  "capabilities": ["route-optimization", "live-eta"]
}
\end{lstlisting}

The SDK provides \texttt{resolveDidWeb()},
\texttt{didDocumentMatchesAddress()}, and
\texttt{buildOperatorDidDocument()} helpers in
\texttt{@figaro/core/extensions}. The metadata document is the
extension point at which ERC-8004 identity descriptors can be
adopted by any deployment that wants them; no kernel-layer
contract change is required.

\paragraph{Why agents are wallet-tier, not entity-tier.}
A consequence of the actor-neutrality property is that the
architecture does not register agents as first-class on-chain
entities. ERC-8004's Identity Registry takes the alternative
route: each agent is a registered entity with its own NFT-backed
handle. Adopting that registry as a kernel-level requirement
would re-introduce the platform-as-vetter problem the actor-
neutrality property is designed to eliminate. A wallet-tier
agent is the consistent extension of kernel-level
actor-neutrality one tier up: the wallet is the agent's identity,
and the agent's reputation is its on-chain settlement history.
ERC-8004 identifiers can travel in the metadata document for
deployments that want them, but they are optional discoverability
metadata, not a registration the architecture requires.

\paragraph{The structural separation.}
Identity descriptors (whether ERC-8004-formatted or otherwise)
are how wallets \emph{find} each other; bonded commitments are
how they \emph{trust} each other. The two are intentionally
separate. A wallet that publishes service descriptors is
discoverable but not yet trusted; a counterparty that wants to
interact reads the descriptors, evaluates the self-published
capabilities, and then bonds against performance. The bond is
the trust layer; the descriptor is the discovery layer.
Conflating the two --- by, for instance, requiring discovery to
imply trust or by requiring trust to imply discovery --- would
re-introduce the platform-as-vetter problem the actor-neutrality
property eliminates. The architecture keeps them separate.

\paragraph{Capability claims are unverified at the kernel layer.}
The \texttt{capabilities} array in an agent's metadata is
\emph{self-asserted}. The kernel does not verify that an agent
who claims \texttt{"route-optimization"} can in fact optimize
routes; the kernel's enforcement is on bonded performance, not
on capability claims. A counterparty that bonds against an agent's
performance is exposed only to bond loss if the agent under-performs;
the counterparty's view of the agent's capability is informed by
the agent's prior on-chain settlement history (which the counterparty
can read directly from the chain) and by any external attestation
the counterparty chooses to consult, but is not certified by any
component of the protocol. This is the appropriate posture for
permissionless infrastructure: capability claims are signals;
performance is enforcement.

% ════════════════════════════════════════════════════════════════════
\section{Scope and Security Posture}
\label{sec:scope}

The architecture is bounded by several scope conditions and operational
constraints that any deployment must respect.

\paragraph{Hardware-isolated signing for production.}
A long-running agent has wallet access by design. Hot-key compromise
is catastrophic in proportion to the agent's bond posture. Production
deployments use hardware-isolated signers (HSMs, secure enclaves,
multi-sig with hardware-key co-signers) rather than software keys
loaded into the agent's runtime. This is standard guidance for any
high-value autonomous wallet, but the bonded architecture amplifies
its importance because the agent's exposure can be substantial.

\paragraph{Default-HITL for high-value actions.}
The architecture's default policy is human-in-the-loop. Autonomous
mode is for actions whose worst case is bounded by predictable
parameters (a single auction claim under \$X; a routine attestation
that reflects an externally-observed fact). For commits, resolutions,
and any action affecting other parties' bonds, the recommended
posture is HITL until the operator has accumulated sufficient
operational confidence in the autonomous policy. Migration from HITL
to autonomous is a deployment decision that follows from observed
policy behavior across many decisions, not a default the architecture
recommends.

\paragraph{Defense-in-depth in the policy.}
The proposer is correct by construction; the policy should still
defend against unexpected proposer behavior as a defense-in-depth
layer. Re-derive bond amounts from process state in the policy. Re-check
counterparty addresses against expected sets. Bound gas in the
execute layer. The policy is the operator's last line of defense
against any architectural assumption that fails in operational
context.

\paragraph{Logging and observability.}
Settlement disputes off-chain need an audit trail; the chain has
half the story. An operational deployment logs proposed actions,
policy decisions, executed transactions, and any divergence between
the agent's expected and observed state. Off-chain forums
adjudicating bonded commitments treat the on-chain settlement
record as evidentiary input; the agent's audit trail is the
operational complement on the agent's side. This is
standard operational practice for any autonomous system; the
bonded architecture invites it because the agent's actions have
material consequences that may be reviewed off-chain.

\paragraph{What this is not.}
The architecture is not a coordination strategy. The factotum
executes what the proposer suggests within the policy's authorization;
it does not decide which markets to enter, which prices are
profitable, which counterparties to trust at the discovery layer,
or which strategic objectives to pursue. The strategy lives in
the policy and in the deployment's higher-layer orchestration. The
architecture is the substrate; the strategy is the operator's.

The architecture is also not LLM-specific. The policy interface
admits any decision-logic. LLM-augmented policies are one
specialization; rule-based policies, reinforcement-learning
policies, model-predictive-control policies, and human-driven
policies are equally well supported. The architecture's
LLM-friendliness derives from the policy interface's generality,
not from any LLM-specific design.

% ════════════════════════════════════════════════════════════════════
\section{Conclusion}
\label{sec:conclusion}

Coordination among mutually-untrusted counterparties has been hard
because the enforcement layer has had nowhere to live without a
centralized institution; the same gap blocks coordination among
autonomous agents, where the institutional vetter is structurally
absent. The bonded settlement primitive's actor-neutrality property
turns bonded commitments into the missing enforcement layer: the
same equilibrium argument that makes cooperation weakly dominant
for any participant makes it weakly dominant regardless of whether
the participant is a human, an autonomous agent, or a hybrid system.
The bond posture is type-blind by construction.

The Agent SDK presented in this paper is a reference architecture
for wallet-bearing clients: a stateful loop with sync, propose,
policy, and execute phases, structured as a control-theory
controller-on-plant architecture. The SDK is convenience tooling,
not a protocol requirement. Reference policies cover typical
operational roles (HITL operator, buyer-with-budget,
seller-of-record, Dutch-auction bidder, auditor) and apply
identically when the wallet is held by a human approving each
action through a UI, by an autonomous agent running its own
decision logic, or by a hybrid system. LLM-augmented policies are
a natural specialization for non-trivial decision logic, with the
bonded architecture's per-action exposure bound and the policy's
aggregate-cap mechanism turning LLM-policy risk into a decomposable
and analyzable surface. Discoverability via ERC-8004 identity
descriptors and \texttt{did:web} is metadata, not protocol; the
discovery layer and the trust layer are separate by design, and
ERC-8004's reputation and validation registries are subsumed by
the bonded primitive at the kernel layer.

The contribution is architectural rather than novel mechanism
design: actor-neutrality is the load-bearing claim, and the SDK is
the operational artifact that demonstrates the claim is implementable.
The participants are not protocol-aware in any privileged sense;
they sign, they bond, they perform. The substrate makes their
participation legible to their counterparties and economically
rational for them to undertake. Agent-only coordination is the
extreme case the architecture admits; coordination between humans,
between agents, and across the human-agent boundary is the rule.

\section*{Acknowledgments}
The Agent SDK and the factotum reference were developed alongside
the kernel and the runtime; they are operational realizations of the
actor-neutrality property the bonded settlement primitive carries
by construction.

% ════════════════════════════════════════════════════════════════════
\begin{thebibliography}{99}

\bibitem[Shoham and Leyton-Brown(2009)]{shoham_leyton_brown_2009}
Shoham, Y. and Leyton-Brown, K.
\newblock \emph{Multiagent Systems: Algorithmic, Game-Theoretic, and
  Logical Foundations}.
\newblock Cambridge University Press, Cambridge, 2009.

\bibitem[Parkes and Seuken(2024)]{parkes_seuken_2024}
Parkes, D.~C. and Seuken, S.
\newblock \emph{Economics and Computation: An Introduction to
  Algorithmic Game Theory, Computational Social Choice, and Fair
  Division}, 2nd edition.
\newblock Cambridge University Press, Cambridge, 2024.

\bibitem[Milgrom(2017)]{milgrom_2017}
Milgrom, P.
\newblock \emph{Discovering Prices: Auction Design in Markets with
  Complex Constraints}.
\newblock Columbia University Press, New York, 2017.

\bibitem[Khalil(2002)]{khalil_2002}
Khalil, H.~K.
\newblock \emph{Nonlinear Systems}, 3rd edition.
\newblock Prentice Hall, Upper Saddle River, NJ, 2002.

\bibitem[Bertsekas(2005)]{bertsekas_2005}
Bertsekas, D.~P.
\newblock \emph{Dynamic Programming and Optimal Control}, 3rd edition.
\newblock Athena Scientific, Belmont, MA, 2005.

\bibitem[W3C(2022)]{w3c_did_2022}
W3C.
\newblock Decentralized Identifiers (DIDs) v1.0: Core Architecture,
  Data Model, and Representations.
\newblock W3C Recommendation, July 2022.

\bibitem[ERC-8004(2026)]{erc8004}
De~Rossi, M., Crapis, D., Ellis, J., and Reppel, E.
\newblock {ERC-8004}: Trustless Agents.
\newblock Ethereum Improvement Proposal, finalized 2026.
\newblock \url{https://eips.ethereum.org/EIPS/eip-8004}

\bibitem[Yao et~al.(2023)]{yao2023react}
Yao, S., Zhao, J., Yu, D., Du, N., Shafran, I., Narasimhan, K.,
  and Cao, Y.
\newblock {ReAct}: Synergizing Reasoning and Acting in Language
  Models.
\newblock In \emph{Proc.\ International Conference on Learning
  Representations (ICLR)}, 2023.

\bibitem[Shinn et~al.(2023)]{shinn2023reflexion}
Shinn, N., Cassano, F., Berman, E., Gopinath, A., Narasimhan, K.,
  and Yao, S.
\newblock Reflexion: Language Agents with Verbal Reinforcement
  Learning.
\newblock In \emph{Proc.\ Advances in Neural Information
  Processing Systems (NeurIPS)}, 2023.

\bibitem[Park et~al.(2023)]{park2023generative}
Park, J.~S., O'Brien, J.~C., Cai, C.~J., Morris, M.~R., Liang,
  P., and Bernstein, M.~S.
\newblock Generative Agents: Interactive Simulacra of Human
  Behavior.
\newblock In \emph{Proc.\ ACM Symposium on User Interface
  Software and Technology (UIST)}, 2023.

\bibitem[Wu et~al.(2023)]{wu2023autogen}
Wu, Q., Bansal, G., Zhang, J., Wu, Y., Li, B., Zhu, E., Jiang, L.,
  Zhang, X., Zhang, S., Liu, J., Awadallah, A.~H., White, R.~W.,
  Burger, D., and Wang, C.
\newblock {AutoGen}: Enabling Next-Gen {LLM} Applications via
  Multi-Agent Conversation.
\newblock arXiv preprint, 2023.

\bibitem[Greshake et~al.(2023)]{greshake2023injection}
Greshake, K., Abdelnabi, S., Mishra, S., Endres, C., Holz, T.,
  and Fritz, M.
\newblock Not what you've signed up for: Compromising Real-World
  {LLM}-Integrated Applications with Indirect Prompt Injection.
\newblock In \emph{Proc.\ ACM Workshop on Artificial Intelligence
  and Security (AISec)}, 2023.

\bibitem[Dong et~al.(2024)]{dong2024agentops}
Dong, L., Lu, Q., and Zhu, L.
\newblock {AgentOps}: Enabling Observability of {LLM} Agents.
\newblock arXiv preprint arXiv:2411.05285, 2024.
\newblock \url{https://arxiv.org/abs/2411.05285}

\end{thebibliography}

\end{document}
